\section{Conclusion}
\label{sec:conclusion}

We began by observing that the verbosity preference of RLHF annotators in HH-RLHF flips sign across annotation rounds --- from penalizing longer responses to rewarding them --- tracking the stylistic profile of the AI text they were evaluating. Our experiments provide the first computational evidence that this pattern is a directional shift specifically aligned with the AI-typicality axis in sentence embedding space.

Our main contributions are: (1) the first coefficient-resolution measurement of verbosity preference reversal across annotation strata ($\Delta\beta_L = +0.080$, non-overlapping 95\% CI); (2) discriminant-valid AI-typicality geometric projection ($\beta_\text{exposure} = 0.041$, $p = 2.05\times10^{-5}$; placebo $p = 0.48$), validating AAI components 1--2; and (3) a concrete data requirements specification for confirming individual annotator causal adaptation.

The highest-priority follow-on is completing H-M3: testing whether the verbosity upweighting in later-round labels propagates to reward model scoring behavior ($\sim$4--8 GPU hours). Purpose-built annotation experiments with genuine within-annotator tracking would enable the causal attribution test that distinguishes individual adaptation from cohort composition effects.

RLHF pipelines are built on the assumption that human annotators provide a stable signal of what constitutes a good response. Our findings suggest this assumption deserves empirical scrutiny in any dataset with multi-round annotation structure. The alignment process that is supposed to make AI systems more human may, over annotation rounds, also make human annotators more AI-typical. Measuring this dynamic --- with instruments like the AAI --- is a necessary step toward annotation pipelines that remain anchored to genuine human values rather than the stylistic artifacts of the systems they are meant to evaluate.
